% =====================================================================
%  Conclusion generale
% =====================================================================
\cleardoublepage
\phantomsection
\addcontentsline{toc}{chapter}{Conclusion generale}
\chapter*{Conclusion generale}
\markboth{Conclusion generale}{}

\section*{Rappel de la demarche}
\addcontentsline{toc}{section}{Rappel de la demarche}

Ce projet de fin d'etudes partait d'une question posee en introduction :
comment doter une entreprise de transport de marchandises d'un systeme
d'information unifie qui centralise l'exploitation de sa flotte,
industrialise sa maintenance, ouvre un canal numerique a ses clients et
mette l'aide a la decision au service de ses choix quotidiens ?

L'etude prealable a montre que les pratiques en place, feuilles de
calcul, echanges informels et tableau de bord alimente a la main,
atteignaient leur limite des que plusieurs personnes devaient consulter
ou modifier la meme information. Aucune des trois familles d'outils
examinees ne couvrait l'ensemble du besoin, les manques les plus nets
portant sur l'aide a la decision et sur la surveillance des echeances
reglementaires.

L'analyse a traduit ce constat en exigences, puis le travail a ete
conduit selon le cadre Scrum, en six sprints d'un mois environ, chacun
livrant un domaine metier utilisable de bout en bout.

\section*{Resultats obtenus}
\addcontentsline{toc}{section}{Resultats obtenus}

Le perimetre realise comprend \essentiel{huit services metier} en .NET~8
appliquant la Clean Architecture, \essentiel{quatre services d'aide a la
decision} en Python, \essentiel{trois applications} destinees a des
publics distincts, et une infrastructure de vingt conteneurs orchestres
par Docker Compose autour d'une base relationnelle, d'un courtier de
messages et d'une base vectorielle. Les six sprints totalisent
\essentiel{trois cent six points de carnet, integralement livres}.

Les quatre objectifs fixes au depart ont ete atteints. Le referentiel est
\essentiel{unifie} : vehicules, equipements, personnels et clients sont
decrits une seule fois. La maintenance est \essentiel{industrialisee} :
ordres de travail suivis par etat, interventions tracees, stock surveille.
Le client dispose d'un \essentiel{canal autonome}, de la demande au suivi
en direct. Enfin, l'aide a la decision est \essentiel{integree au flux de
travail} plutot que releguee dans un module separe.

\section*{Difficultes rencontrees}
\addcontentsline{toc}{section}{Difficultes rencontrees}

Trois difficultes ont marque le deroulement du projet.

\paragraph{Le decoupage en services.}
Determiner ou passe la frontiere entre deux services est plus delicat que
de la coder. La question s'est posee des le sprint~3 : la demande de
transport releve-t-elle du domaine des missions ou de celui de la
facturation ? La reponse retenue, la rattacher aux missions, decoule de
l'observation que son cycle de vie s'acheve a la creation de la mission,
et non au reglement.

\paragraph{La coherence du vocabulaire entre services.}
Chaque service definit ses propres enumerations. L'un d'eux, le service
d'alertes, a fait diverger la sienne de celle du projet commun. Ecrire une
valeur du vocabulaire commun qui n'y figurait pas suffisait a faire
echouer la lecture. \essentiel{Cette difficulte est structurelle} : elle
est le revers de l'autonomie des services.

\paragraph{La restitution graphique de la modelisation.}
Produire des diagrammes UML lisibles s'est revele un travail a part
entiere. Certains atteignaient onze mille pixels de large, d'autres
presentaient des associations contournant la frontiere du systeme. Il a
fallu etablir des regles de mise en page explicites, puis les appliquer
de maniere systematique, pour que les modeles soient exploitables dans un
document imprime.

\section*{Apports personnels}
\addcontentsline{toc}{section}{Apports personnels}

Sur le plan \essentiel{technique}, ce projet a permis de mettre en oeuvre
une architecture microservices complete, de la conception au deploiement
conteneurise, avec deux piles technologiques et trois types
d'applications clientes. La Clean Architecture, appliquee huit fois, a
cesse d'etre une notion pour devenir une habitude de travail.

Sur le plan \essentiel{methodologique}, la conduite en sprints a montre
sa valeur au moment ou le perimetre a evolue : plusieurs fonctionnalites,
dont l'archivage reversible des comptes et la maintenance predictive, ont
ete identifiees en cours de projet. Un cycle sequentiel les aurait
renvoyees a une version ulterieure.

Sur le plan \essentiel{professionnel}, l'exercice le plus formateur a ete
d'apprendre a \essentiel{documenter ce qui ne fonctionne pas comme prevu}
plutot qu'a le passer sous silence. Les ecarts par rapport au patron de
reference, base partagee, absence de cache et de chiffrement, sont
consignes dans les diagrammes eux-memes. C'est une pratique que
l'entreprise d'accueil, \entreprise, encourage, et qui rend une
architecture reprenable par un tiers.

\section*{Perspectives}
\addcontentsline{toc}{section}{Perspectives}

Le travail realise ouvre plusieurs prolongements, classes ici par ordre
de priorite.

\begin{table}[H]
\centering
\caption{Perspectives d'evolution}
\label{tab:perspectives}
\small
\begin{tabularx}{\textwidth}{@{}p{3.4cm}X@{}}
\toprule
\textbf{Axe} & \textbf{Contenu} \\ \midrule
Mise en securite & Terminaison du chiffrement a la passerelle et
  redirection du trafic en clair : prealable a toute mise en ligne \\
Exploitation du bus & Sur vingt-deux evenements declares, un seul est
  publie et un seul consomme ; les exploiter reduirait les lectures
  directes entre domaines \\
Cloisonnement des donnees & Passer a une base par service, en remplacant
  les lectures en SQL direct par des appels ou des evenements \\
Bascule vers l'apprentissage & Substituer un modele appris au moteur de
  regles des qu'un historique suffisant sera constitue \\
Mise a l'echelle & Cache applicatif et orchestrateur de grappe \\
Application chauffeur & Mise en file locale des positions lorsque la
  couverture reseau fait defaut \\
\bottomrule
\end{tabularx}
\end{table}

\section*{Mot de la fin}
\addcontentsline{toc}{section}{Mot de la fin}

La plateforme livree n'est pas un prototype : elle fonctionne, elle est
deployee et elle repond aux besoins exprimes. Elle n'est pas non plus un
produit acheve, et le present rapport s'est attache a dire ou passe cette
limite. Cette distinction est peut-etre l'apprentissage principal de ces
\stageDuree{} mois : un systeme d'information se juge autant a ce qu'il
assume ne pas faire qu'a ce qu'il realise.
